\documentclass[letter]{article}
\usepackage{color}
\usepackage{graphicx}
\usepackage{pifont}
\usepackage[top=2in, bottom=1.5in, left=1in, right=1in]{geometry}
\usepackage{array,booktabs}

\newcolumntype{M}[1]{>{\centering\arraybackslash}m{#1}}

\setlength{\parskip}{0.3cm plus4mm minus3mm}
\setlength{\parindent}{0cm}

\newcommand{\clubs}{{\large\Pisymbol{psy}{"A7}}}
\newcommand{\spades}{{\large\Pisymbol{psy}{"AA}}}
\newcommand{\hearts}{{\large\color{red}\Pisymbol{psy}{"A9}}}
\newcommand{\diamonds}{{\large\color{red}\Pisymbol{psy}{"A8}}}

\newcommand{\clubsX}[1]{\clubs\,#1}
\newcommand{\spadesX}[1]{\spades\,#1}
\newcommand{\heartsX}[1]{\hearts\,#1}
\newcommand{\diamondsX}[1]{\diamonds\,#1}

\newcommand{\trumpjacks}{\clubsX{J} \spadesX{J} \heartsX{J} \diamondsX{J}}
\newcommand{\nonjackcards}{{\bfseries A 10 K Q 9 8 7}}
\newcommand{\cards}{{\bfseries A K Q J 10 9 8 7}}
\newcommand{\bids}{{\bfseries18, 20, 22, 23, 24, 27, 30, 33, 35, 36, 40, 44,
    45, 46, 48, 50, 54, 55, 59, 60} etc.}


% Most content from:
% http://www.pagat.com/schafk/skat.html#gamevalue

\begin{document}

  \section*{Overview}

  Skat is a three-handed trick taking game. It can also be played by four
  people with the dealer sitting out. Each player gets 10 of the 32 cards and
  the remaining two form the \emph{skat} and are placed face down in the middle
  of the table.

  The game starts with an auction of which the winner becomes the
  \emph{declarer}. This player has to play alone against the other two, who play
  as a team. The declarer has the right to exchange the skat cards for two of
  the own cards. The declarer also chooses the trump suit or instead to play
  \emph{Grand} (jacks are the only tump) or \emph{Null} (no trumps and declarer
  must not get any tricks).

  All \emph{card points} combined add up to 120, while some cards don't
  contribute any points. To win, declarer has to to take at least 61 points. For
  the openents, 60 points is enough to win.

  Each game has a \emph{value}, in \emph{game points}, that depends on the
  trumps or game chosen, the location of the top trumps (\emph{matadors)} and
  whether the declarer used the skat. The declarer generally wins the value of
  the game if successful, and loses the game value (doubled if skat was
  exchanged) if unsuccesful. The game points one wins and loses are separate
  from the card points (which determine whether one wins or loses)


  \section*{Cards}
  
  Skat is often played with German suits. However, French suited cards are
  common as well. There are 32 cards in use. For each suit there is: A K Q J 10
  9 8 7:

  \begin{tabular}{ M{0.5cm} m{2cm} M{0.5cm} m{2cm} m{2cm} }
    {\huge \clubs} & clubs & \includegraphics[height=0.5cm]{images/german_suit_acorn} & acorns \\
    {\huge \spades} & spades & \includegraphics[height=0.5cm]{images/german_suit_gras} & leaves \\
    {\huge \hearts} & hears & \includegraphics[height=0.5cm]{images/german_suit_herz} & herts \\
    {\huge \diamonds} & diamonds & \includegraphics[height=0.5cm]{images/german_suit_schell} & bells \\
    {\bfseries \large A} & ace & {\bfseries \large A} & ace \\
    {\bfseries \large K} & king & {\bfseries \large K} & king \\
    {\bfseries \large Q} & queen & {\bfseries \large O} & ober \\
    {\bfseries \large J} & jack & {\bfseries \large U} & unter \\
  \end{tabular}

  \pagebreak

  \subsection*{Ranking of cards and games}
 
  The ranking depends on the game the declarer chooses to play. 

  \subsubsection*{Suit games}

  All four jacks are trump, regardless of the suit. Their ranking follows the
  ranking of suits: \clubs \spades \hearts \diamonds. Next and also trump are the
  remaining seven cards of the chosen suit. This makes eleven trumps in total:
  \trumpjacks \nonjackcards. The other three suits contain just seven cards,
  ranked: \nonjackcards.

  \subsubsection*{Grand}

  The four jacks are the only tump: \trumpjacks. The remaining four suits each
  contain seven cards ranking from high to low: \nonjackcards.

  \subsubsection*{Null}

  There are no trumps. The eight cards of each suit rank from high to low:
  \cards. Notice that they are \emph{not} ordered by card value.

  \section*{Deal}

  The first dealer is chosen at random. Thereafter the turn to deal rotates
  clockwise. The dealer shuffles an the player to the dealer's right cuts.
  Cards are dealt face down in four steps: 1. Three cards to each player 2. Two
  cards in the center (the skat) 3. Four cards to each player and finally 4.
  Three cards to each player.

  \subsection*{Card values}

  In suit games and grand, the cards have the following values:

  {\bfseries
    \begin{tabular}{c c c c c c c c c c}
      J &  A & 10 & K & Q & 9 & 8 & 7 \\
      2 & 11 & 10 & 4 & 3 & 0 & 0 & 0 \\
    \end{tabular}
  }

  \section*{Auction}

  Each bid is a number which is the \emph{game value} in \emph{game points} of
  some possivle game (see below for calculation of game values). The possible
  bids are therefore \bids. If you bid or accept a bid it meansyou are prepared
  to play a contract of at least that value in game points.

  The player to the dealer's left is called \emph{forehand (F)}, the player to
  the forehand's left \emph{middlehand (M)}, and the player to the middlehand's
  left is \emph{rearhand (R)}. If there are three players, R is the dealer. If
  there are four players, R is to the dealer's right. Throughout bidding F is
  senior to M wo is senior to R. The principle is that a senior player only has
  to \emph{equal} a junior player's bid to win the auction, whereas a junior
  player has to bid \emph{higher} than a senior player to win.

  The first part of the auction takes place between F and M. M speaks first,
  either passing or bidding a number. There is no advantage in making a higher
  than necessary bid so M will normally either pass or begin with the lowest
  bid: 18. If M bids a number, F can either give up the chance to be declarer by
  saying "pass" or compete by saying "yes", which means that F bids the same
  number that M just bid. If F says "yes", M can say "pass", or continue the
  auction with a higher bid, to which F will again answer "yes" or "pass".
  This continues until either F or M drops out of the auction by passing - once
  having passed you get no further opportunity to bid on that hand. 

  The second part of the auction is similar to the first part, but takes place
  between R and the survivor of the first part (i.e. whichever of F and M did
  not pass. As the junior player, R either passes or bids a succession of
  numbers, the first of which must be higher than any number mentioned in the
  first part of the auction. To each number bid by R, the survivor must answer
  "yes" or "pass". The winner of the second part of the auction becomes the
  declarer, and the bid is the last number the declarer said or accepted.

  If both M and R pass without having bid, then F can either be declarer at the
  lowest bid (18), or can throw in the cards without play. If the cards are
  thrown in there is no score for the hand, and the next dealer deals.

  Examples:

  \begin{minipage}[t][5cm]{0.5\textwidth}
    \centering
    \vspace{0pt}
    \begin{tabular}{l l l l}
         F &    M &  R & \\
       \cmidrule(r){1-3}
           &   18 &    & \\
       yes &   20 &    & \\
       yes & pass &    & (F wins first part) \\
           &      & 22 & \\
       yes &      & 23 & \\
       yes &      & 24 & \\
      pass &      &    & (R is declarer in 24) \\
    \end{tabular}
    \vfill
  \end{minipage}
  \begin{minipage}[t][5cm]{0.5\textwidth}
    \centering
    \vspace{0pt}
    \begin{tabular}{l l l l}
        F &    M &  R & \\
       \cmidrule(r){1-3}
          & pass &    & (F wins first part) \\
          &      & 18 & \\
      yes &      & pass & (F is declarer in 18) \\
    \end{tabular}
    \vfill
    \begin{tabular}{l l l l}
         F &    M &  R & \\
       \cmidrule(r){1-3}
           &   18 &    & \\
      pass &      &    & (M wins first part) \\
           &      & 20 & \\
           & pass &    & (R is declarer in 20) \\
    \end{tabular}
  \end{minipage}

  To remember whose turn it is to start the bidding, German players sometimes
  say "geben, hören, sagen" (deal, listen, speak), pointing in turn to dealer,
  forehand and middlehand. If middlehand forgets to begin, forehand can start
  proceedings by saying "I'm forehand" or "I'm listening", or "Speak to me!".

  \section*{Possible Contracts}

  \section*{The Play}

  \section*{Calculating the Value of the Game}

  \section*{Scoring}
\end{document}
